% Options for packages loaded elsewhere
\PassOptionsToPackage{unicode}{hyperref}
\PassOptionsToPackage{hyphens}{url}
\documentclass[
  ignorenonframetext,
]{beamer}
\newif\ifbibliography
\usepackage{pgfpages}
\setbeamertemplate{caption}[numbered]
\setbeamertemplate{caption label separator}{: }
\setbeamercolor{caption name}{fg=normal text.fg}
\beamertemplatenavigationsymbolsempty
% remove section numbering
\setbeamertemplate{part page}{
  \centering
  \begin{beamercolorbox}[sep=16pt,center]{part title}
    \usebeamerfont{part title}\insertpart\par
  \end{beamercolorbox}
}
\setbeamertemplate{section page}{
  \centering
  \begin{beamercolorbox}[sep=12pt,center]{section title}
    \usebeamerfont{section title}\insertsection\par
  \end{beamercolorbox}
}
\setbeamertemplate{subsection page}{
  \centering
  \begin{beamercolorbox}[sep=8pt,center]{subsection title}
    \usebeamerfont{subsection title}\insertsubsection\par
  \end{beamercolorbox}
}
% Prevent slide breaks in the middle of a paragraph
\widowpenalties 1 10000
\raggedbottom
\AtBeginPart{
  \frame{\partpage}
}
\AtBeginSection{
  \ifbibliography
  \else
    \frame{\sectionpage}
  \fi
}
\AtBeginSubsection{
  \frame{\subsectionpage}
}
\usepackage{iftex}
\ifPDFTeX
  \usepackage[T1]{fontenc}
  \usepackage[utf8]{inputenc}
  \usepackage{textcomp} % provide euro and other symbols
\else % if luatex or xetex
  \usepackage{unicode-math} % this also loads fontspec
  \defaultfontfeatures{Scale=MatchLowercase}
  \defaultfontfeatures[\rmfamily]{Ligatures=TeX,Scale=1}
\fi
\usepackage{lmodern}
\ifPDFTeX\else
  % xetex/luatex font selection
\fi
% Use upquote if available, for straight quotes in verbatim environments
\IfFileExists{upquote.sty}{\usepackage{upquote}}{}
\IfFileExists{microtype.sty}{% use microtype if available
  \usepackage[]{microtype}
  \UseMicrotypeSet[protrusion]{basicmath} % disable protrusion for tt fonts
}{}
\makeatletter
\@ifundefined{KOMAClassName}{% if non-KOMA class
  \IfFileExists{parskip.sty}{%
    \usepackage{parskip}
  }{% else
    \setlength{\parindent}{0pt}
    \setlength{\parskip}{6pt plus 2pt minus 1pt}}
}{% if KOMA class
  \KOMAoptions{parskip=half}}
\makeatother
\setlength{\emergencystretch}{3em} % prevent overfull lines
\providecommand{\tightlist}{%
  \setlength{\itemsep}{0pt}\setlength{\parskip}{0pt}}
% Slide layout defaults for warning-clean scientific decks.
\usepackage{etoolbox}
\IfFileExists{xurl.sty}{\usepackage{xurl}}{}
\IfFileExists{seqsplit.sty}{\usepackage{seqsplit}}{\newcommand{\seqsplit}[1]{#1}}
\protected\def\breakseq#1{\seqsplit{#1}}
\protected\def\breaktt#1{\begingroup\ttfamily\seqsplit{#1}\endgroup}
\setlength{\emergencystretch}{6em}
\tolerance=5000
\hbadness=10000
\hfuzz=1pt
\setlength{\tabcolsep}{2pt}
\AtBeginEnvironment{longtable}{\tiny\renewcommand{\arraystretch}{0.86}\setlength{\tabcolsep}{1pt}}
\AtBeginEnvironment{tabular}{\tiny\renewcommand{\arraystretch}{0.86}\setlength{\tabcolsep}{1pt}}

% Natbib and cross-reference fallbacks — slides don't load natbib
% or cleveref, but combined-PDF manuscript prose may emit these
% commands. The fallback renders citations as a bracketed key list
% and cross-references as detokenized labels so slides don't fail on
% undefined control sequences or raw underscores. \providecommand is
% a no-op if packages load later.
\providecommand{\citep}[1]{[#1]}
\providecommand{\citet}[1]{#1}
\providecommand{\citealp}[1]{#1}
\providecommand{\citeauthor}[1]{#1}
\providecommand{\citeyear}[1]{#1}
\providecommand{\cref}[1]{\texttt{\detokenize{#1}}}
\providecommand{\Cref}[1]{\texttt{\detokenize{#1}}}

\newtheorem{proposition}{Proposition}
\newtheorem{hypothesis}{Hypothesis}
\usepackage{bookmark}
\IfFileExists{xurl.sty}{\usepackage{xurl}}{} % add URL line breaks if available
\urlstyle{same}
% fallback for those not using the hyperref driver hyperxmp:
\makeatletter
\@ifundefined{xmpquote}{\newcommand{\xmpquote}[1]{#1}}{}
\makeatother
\hypersetup{
  hidelinks,
  pdfcreator={LaTeX via pandoc}}

\author{\texorpdfstring{}{}}
\date{}

\begin{document}

\section{Modeling Lens: Active Inference Without Mechanism
Claims}\label{sec:active_inference}

\begin{frame}[fragile,allowframebreaks]{Modeling Lens: Active Inference
Without Mechanism Claims}
Active inference gives DigiPPPiP a disciplined modeling vocabulary for
dyadic mark-making. Each partner maintains expectations about the shared
canvas, the other partner's contribution, and the evolving meaning of
the artifact. The shared drawing is therefore a sensory sample, a
prediction target, and an action channel \breakseq{[@friston2010fep;
@friston2017process; @ramstead2020two]}. Active-inference accounts of
human communication support this interactional reading only when the
model stays tied to observable communicative events rather than becoming
a generic metaphor \breakseq{[@vasil2020communication]}. In second-person
terms, the partner is not a stimulus to observe from outside but an
agent whose responses shape one's own policy choices
\breakseq{[@schilbach2013secondperson; @redcay2019secondperson;
@bolis2024secondperson]}.

The main claim is intentionally narrow. DigiPPPiP does not claim that
participants literally compute free-energy equations, nor that the
current simulations measure a real couple. The claim is that a useful
formal account must specify states, observations, policies, update
rules, and observable consequences. If a study cannot map interface
events to observations, define partner-relevant policies, or distinguish
reciprocal updating from shared-stimulus effects, then active inference
is not doing explanatory work for that study.

The formal appendix consolidates the mathematical primitives, model
figures, and toy diagnostics in one place
\breakseq{[@sec:formalisms\_appendix]}. That appendix maps DigiPPPiP events to
active-inference variables, gives the Gaussian point-belief model used
by \breaktt{src/active\_inference.py}, shows the coupled-versus-decoupled
toy run, and places narrative entropy, epistemic arcs, and geometric
hyperscanning under the same evidence boundary. Keeping those details
out of the main line makes the argument easier to inspect: the
manuscript first states why the model is useful, then the appendix shows
exactly what the model assumes.

This separation also protects the empirical agenda. Hyperscanning,
narrative information, and free-energy trajectories are candidate
measurement languages, not outcome evidence. Synchrony can be inflated
by shared stimulus timing, motion artifacts, physiology, preprocessing
decisions, and analytic windows \breakseq{[@tachtsidis2016fnirs;
@hamilton2021hyperscanning; @zimmermann2024methodological]}. Narrative
entropy can identify symbolic structure without proving shared meaning.
A conceptual free-energy curve can expose a coupling assumption without
showing that partners improved. The framework is strongest when each
formalism names the future comparison that could reject it.

The practical upshot is a modeling checklist. A DigiPPPiP study should
log the visible mark, the timing of the response, the partner-facing
cue, the repair or withdrawal option, and the artifact-governance state.
It should then ask whether the model explains something beyond simpler
alternatives such as ordinary shared drawing, video conversation, or
solo art-making. Active inference remains in the manuscript because it
forces those commitments into the open, not because it upgrades
conceptual simulations into evidence.
\end{frame}

\end{document}
